\documentclass[12pt]{article}
\usepackage{latexsym, amscd, %graphicx, 
amssymb}
\newcommand{\SECT}[1]
		{\vspace{.3cm}
		$\!\!\!\!\!\!\!\!\!\!\!${\bf {#1}}
		\vspace{.2cm}}

\begin{document}


\begin{center}

{\large \bf Math 152, Fall, 2004, Workshop 4\par}
\vspace{.5cm}
{{\bf Honors Section} 
\par}
\end{center}

\begin{enumerate}

\item Calculate each of the indefinite integrals (these are all pretty easy): 

\begin{enumerate}

\item ${\displaystyle \int \sec^{4}x dx}.$

\item ${\displaystyle \int \sec^{3} x \tan x dx}$.

\item ${\displaystyle \int \sec^{2}x \tan^{2}x dx}$.

\item ${\displaystyle \int \sec x \tan^{3} x dx}$.

\end{enumerate}

\item Calculate the area which lies inside both of the ellipses 
${\displaystyle \frac{x^{2}}{3} + y^{2} = 1}$ and
${\displaystyle  x^{2} + \frac{y^{2}}{3} = 1}$

\item The region $R$ is bounded below by the $x$-axis, bounded on 
the left by the line $x = 0$, bounded 
on the right by the line $x = 2$, and bounded above by the curve 
$$y = \frac{2x + 1}{x^2 + 3x + 2}.$$

\begin{enumerate}

\item Sketch the region $R$ and set up a definite integral that 
gives the area of $R$. Then calculate the 
integral using the method of partial fractions. 

\item The region $R$ is rotated around the $x$-axis to generate 
a solid body $B$. Sketch $B$ and, by slicing 
$B$ into discs, set up a definite integral that gives the 
volume of B. Calculate the integral using
the method of partial fractions.
(warning: there are four undetermined coefficients). 

\end{enumerate}

4. Many integrals can be done with ``rationalizing substitutions" 
which change the integrals into 
integrals involving rational functions. In turn, these integrals 
can be computed (at least theoretically) using partial fractions, 
although often it is simpler to integrate directly or to use a 
trigonometric
substitution. Integrate: 
$$\int \frac{1}{x+\sqrt{x}}dx$$
(try $x = t^{2}$) and  
$$\int \frac{1}{e^{2x} + 1} dx$$
(try $t = e^{x}$). 

\end{enumerate}



\end{document}